% Artigo do projeto de extensão: priorização e roteirização de empilhadeiras com grafos dinâmicos.
% Baseado no modelo abnTeX2 (ABNT NBR 6022:2003).
% Compilar: pdflatex, bibtex, pdflatex, pdflatex (ver latex/README.md).

\documentclass[
	article,
	11pt,
	oneside,
	a4paper,
	english,
	brazil,
	sumario=tradicional
	]{abntex2}

% ---
% Pacotes
% ---
\usepackage{lmodern}
\usepackage[T1]{fontenc}
\usepackage[utf8]{inputenc}
\usepackage{indentfirst}
\usepackage{color}
\usepackage{graphicx}
\usepackage{microtype}
\usepackage{booktabs}			% Tabelas com linhas horizontais limpas
\usepackage{amsmath}			% Equações
\usepackage[alf]{abntex2cite}	% Citações padrão ABNT

% ---
% Dados da capa
% ---
\titulo{Priorização e roteirização de empilhadeiras em centros logísticos com grafos dinâmicos}
\autor{%
	Márcio Brener Cantuária Santos\thanks{Matrícula 239781. Engenharia da Computação, Fundação Mathias Machline.}
	\and
	Heitor de Medeiros Santos\thanks{Matrícula 239834. Engenharia da Computação, Fundação Mathias Machline.}}
\local{Manaus}
\data{2026}
% Orientador: aparece logo abaixo da capa (ver \maketitle).
\newcommand{\orientador}{Prof. Gabriel de Oliveira Araújo}

% ---
% Aparência do PDF
% ---
\definecolor{blue}{RGB}{41,5,195}
\makeatletter
\hypersetup{
	pdftitle={\@title},
	pdfauthor={Márcio Brener Cantuária Santos; Heitor de Medeiros Santos},
	pdfsubject={Priorização e roteirização de empilhadeiras com grafos dinâmicos},
	pdfcreator={LaTeX with abnTeX2},
	pdfkeywords={grafos dinâmicos}{empilhadeiras}{logística}{simulação},
	colorlinks=true,
	linkcolor=blue,
	citecolor=blue,
	filecolor=magenta,
	urlcolor=blue,
	bookmarksdepth=4
}
\makeatother

% ---
% Margens e espaçamentos
% ---
\setlrmarginsandblock{3cm}{3cm}{*}
\setulmarginsandblock{3cm}{3cm}{*}
\checkandfixthelayout
\setlength{\parindent}{1.3cm}
\setlength{\parskip}{0.2cm}
\SingleSpacing

\begin{document}

\frenchspacing

% ----------------------------------------------------------
% ELEMENTOS PRÉ-TEXTUAIS
% ----------------------------------------------------------
\maketitle

\begin{center}
	Orientador: \orientador
\end{center}

\begin{resumoumacoluna}
 Em centros logísticos, a armazenagem de paletes por empilhadeiras ocorre sem
 priorização das tarefas. Como a empilhadeira precisa parar para posicionar,
 elevar e guardar o palete, o trecho da rua onde ela opera fica bloqueado para
 os demais operadores, o que gera esperas e congestionamento, inclusive para
 materiais de baixo giro. Este trabalho propõe um algoritmo baseado em grafos
 dinâmicos que, a cada evento (nova tarefa, entrada ou saída de empilhadeira em
 uma rua), define a ordem de armazenagem e a rota de cada operador. A prioridade
 de cada tarefa combina o giro do material, a posição de destino e a ocupação
 atual da rua, com pesos ajustáveis, e é limitada por uma regra de precedência
 por profundidade, na qual o material de destino mais fundo é armazenado
 primeiro. Um simulador por passos de tempo compara a ordem de chegada com o
 algoritmo, medindo tempo de espera, distância percorrida e número de
 bloqueios. Trata-se de um protótipo, avaliado apenas por simulação em um
 layout genérico, sem aplicação em um centro logístico real. Os resultados
 preliminares mostram melhora no tempo de ciclo dos materiais de maior giro,
 mas a regra de precedência, isoladamente, não reduziu a espera total por
 bloqueio, o que motiva a previsão de conflitos de saída como próxima etapa.

 \vspace{\onelineskip}

 \noindent
 \textbf{Palavras-chave}: grafos dinâmicos. empilhadeiras. logística.
 priorização de tarefas. simulação.
\end{resumoumacoluna}

% ----------------------------------------------------------
% ELEMENTOS TEXTUAIS
% ----------------------------------------------------------
\textual

% ----------------------------------------------------------
\section{Introdução}
% ----------------------------------------------------------

Em um centro logístico, depois que o palete é montado na área de recebimento,
o operador de empilhadeira o busca e o armazena na posição de estoque
correspondente. Em geral, não existe priorização dessas tarefas de armazenagem:
cada operador atende o que lhe é designado, na ordem em que as tarefas surgem.
Essa etapa integra o conjunto de operações de um armazém, que inclui
recebimento, armazenagem, separação de pedidos e expedição
\cite{gu2007}.

Para armazenar, a empilhadeira precisa se posicionar, elevar o palete, guardá-lo
e liberar a rua. Durante esse tempo, ela ocupa o trecho onde está parada e
impede a passagem dos demais. Como essa situação se repete muitas vezes ao dia
e envolve vários operadores simultaneamente, formam-se esperas e
congestionamento na mesma rua, mesmo quando as tarefas em disputa são de
materiais de baixo giro.

Este trabalho propõe tratar o problema com um grafo dinâmico, recalculado a cada
evento, que devolve a ordem de prioridade das tarefas e a rota de cada
operador. O foco do trabalho é o problema de congestionamento e o algoritmo
que o trata. O que se apresenta é um protótipo, avaliado por simulação, sem
aplicação em operação real.

O restante do texto apresenta o problema (Seção~\ref{sec:problema}), os
objetivos (Seção~\ref{sec:objetivos}), a fundamentação teórica
(Seção~\ref{sec:teoria}), a metodologia (Seção~\ref{sec:metodologia}), os
resultados preliminares (Seção~\ref{sec:resultados}) e as considerações finais,
com limitações e trabalhos futuros (Seção~\ref{sec:consideracoes}).

% ----------------------------------------------------------
\section{Definição do problema}
\label{sec:problema}
% ----------------------------------------------------------

Considera-se um centro logístico com docas de recebimento, uma área de
recebimento (pulmão) e ruas de armazenagem com racks nos dois lados. Cada rua
é acessada por um único lado, ou seja, é um beco sem saída, e pertence a um
depósito, que corresponde a uma categoria de material (por exemplo,
medicamentos, produtos de higiene e cuidados pessoais, grandes volumes).

Cada material pertence a uma categoria e a uma curva ABC, e possui endereço
fixo, composto por depósito, rua, posição e nível. A categoria define o
depósito e a rua; a curva define o nível (curva A no nível mais baixo). Como o
destino de cada tarefa já é conhecido pelo material, o algoritmo não escolhe
endereço: ele decide apenas \emph{em que ordem} cada tarefa é executada e
\emph{por qual rota}.

O bloqueio ocorre da seguinte forma: a empilhadeira parada operando ocupa o
trecho onde está, e quem precisa passar por ele, para entrar ou para sair,
espera. Quanto mais alto o nível do endereço, maior o tempo de elevação e,
portanto, maior o tempo em que a rua fica bloqueada. Duas ou mais operações
podem estar na mesma rua ao mesmo tempo, desde que quem está parado não
bloqueie quem ainda precisa passar. Isso impõe uma ordem natural: se a tarefa
da frente da rua é executada antes da tarefa do fundo, a empilhadeira da frente
trava o acesso ao fundo; se o fundo é executado primeiro, as duas operações
coexistem sem espera.

% ----------------------------------------------------------
\section{Objetivos}
\label{sec:objetivos}
% ----------------------------------------------------------

\subsection{Objetivo geral}

Desenvolver e demonstrar um algoritmo baseado em grafos dinâmicos que defina a
ordem de prioridade e a rota das tarefas de armazenagem de empilhadeiras,
considerando o giro dos materiais e a posição de destino, para reduzir esperas
e bloqueios de corredor em um centro logístico.

\subsection{Objetivos específicos}

\begin{alineas}
	\item modelar um layout genérico de centro logístico (docas, pátio, ruas e
	posições) como um grafo com pesos dinâmicos;
	\item definir um índice de prioridade das tarefas combinando o giro do
	material e a posição de destino no depósito;
	\item implementar o algoritmo de priorização e roteirização para múltiplos
	operadores concorrendo pela mesma rua;
	\item disponibilizar uma demonstração visual, de apoio, que execute o
	algoritmo e mostre a ordem de armazenagem sugerida;
	\item comparar, por simulação, o cenário sem priorização (ordem de chegada)
	com o cenário com o algoritmo, quanto a tempo de espera, distância
	percorrida e número de bloqueios.
\end{alineas}

% ----------------------------------------------------------
\section{Fundamentação teórica}
\label{sec:teoria}
% ----------------------------------------------------------

\subsection{Operações de armazém e política de armazenagem}

\citeonline{gu2007} organizam a literatura sobre operações de armazém em
torno das atividades de recebimento, armazenagem, separação de pedidos e
expedição, e discutem os problemas de decisão de cada uma. A tarefa tratada
neste trabalho está na transição entre o recebimento e a armazenagem: levar o
palete da área de recebimento até o seu endereço de estoque.

A literatura de separação de pedidos é revisada por \citeonline{dekoster2007},
que discutem o projeto e o controle dessa atividade, incluindo o layout do
armazém, as políticas de armazenagem e a roteirização dos operadores. Esses
temas são próximos dos que aparecem aqui, mas aplicados ao deslocamento de
empilhadeiras para guardar paletes, e não à coleta de itens.

A relação entre o giro dos materiais e o desempenho do armazém tem origem no
estudo de \citeonline{hausman1976}, que comparam atribuição aleatória de
endereços, atribuição por giro e atribuição por classes de giro em sistemas
automáticos de armazenagem, e mostram que regras baseadas em giro reduzem o
tempo de deslocamento do equipamento. Neste trabalho, o
giro (curva ABC) já determina o nível do endereço, e por isso ele entra no
algoritmo como critério de \emph{prioridade} da tarefa, e não como regra de
escolha de endereço.

No contexto brasileiro, \citeonline{silva2017} combinam a curva ABC com
simulação computacional em um armazém de centro de distribuição, usando o giro
dos produtos como parâmetro para organizar o depósito. Os autores destacam que
os cenários simulados mostram impactos em cada processo que seriam difíceis de
observar no ambiente real sem perturbar a operação, o mesmo motivo pelo qual
este trabalho recorre à simulação.

\subsection{Bloqueio e congestionamento em corredores}

\citeonline{parikh2009} propõem modelos analíticos para estimar o bloqueio
entre operadores em corredores de separação de pedidos e mostram que a largura
do corredor influencia o quanto os operadores se atrapalham. No problema aqui
estudado o bloqueio tem origem parecida, a presença de um veículo parado que
impede a passagem de outro, mas ocorre em ruas de acesso único, nas quais quem
está no fundo depende de quem está na frente para sair.

\subsection{Caminho mínimo e grafos dinâmicos}

O cálculo de rotas neste trabalho usa o algoritmo de \citeonline{dijkstra1959}
para encontrar o caminho de menor custo entre dois nós de um grafo com arestas
ponderadas. Como as condições do armazém mudam a cada evento, o grafo precisa
ser reavaliado ao longo do tempo.

\citeonline{ramalingam1996} estudam a complexidade dos problemas em grafos
dinâmicos, nos quais a solução deve ser atualizada depois de mudanças no grafo,
em vez de recalculada do zero. Em um caso de roteirização dentro de armazém na
literatura nacional, \citeonline{bonassa2011} buscam a rota de separação manual
que minimiza a distância percorrida, com um sistema de apoio à decisão em
planilha. O foco deles é a rota que minimiza a distância, e não a disputa entre
operadores por trechos de rua. A implementação deste protótipo recalcula a rota
a cada decisão, o que é suficiente para o grafo pequeno usado, e a atualização
incremental fica como possibilidade de evolução.

\subsection{Despacho e roteirização de veículos}

A revisão de \citeonline{leanh2006} sobre veículos guiados automaticamente
(AGV) cobre projeto do caminho, escalonamento, roteirização e resolução de
deadlocks. Empilhadeiras conduzidas por operadores não são AGVs, mas o problema
de decidir qual veículo atende qual tarefa e por qual caminho, evitando que
veículos se travem, é análogo ao que se trata aqui.

\citeonline{qiu2002} classificam os algoritmos de escalonamento e roteirização
de AGVs e distinguem os dois problemas. O algoritmo proposto neste trabalho
trata os dois de forma conjunta: o escalonamento define a \emph{ordem} das
tarefas e a roteirização define o \emph{caminho} de cada operador.

Entre as abordagens clássicas, \citeonline{egbelu1984} caracterizam regras
heurísticas de despacho de AGVs em um ambiente de oficina (\emph{job shop}).
A política de ordem de chegada usada como referência neste trabalho é uma regra
de despacho simples desse tipo, e o índice de prioridade proposto é uma regra
mais elaborada, que incorpora giro, posição e ocupação.

\subsection{Síntese}

Nas fontes revisadas, os temas de armazenagem por giro, bloqueio em corredores,
roteirização em grafos e despacho de veículos aparecem em separado
\cite{gu2007,hausman1976,parikh2009,leanh2006,silva2017,bonassa2011}. Não encontramos nelas o
tratamento conjunto do problema aqui descrito: tarefas de armazenagem em ruas
sem saída, com prioridade que combina giro, profundidade e ocupação atual. Essa
combinação é o que o algoritmo proposto busca explorar. Esta síntese se limita
às fontes consultadas e não pretende ser uma revisão sistemática.

% ----------------------------------------------------------
\section{Metodologia}
\label{sec:metodologia}
% ----------------------------------------------------------

\subsection{Natureza e etapas da pesquisa}

Trata-se de uma pesquisa aplicada, de abordagem quantitativa, baseada em
modelagem e simulação. O algoritmo e o simulador foram implementados como
protótipo em JavaScript, sem dependências externas. A simulação foi escolhida
porque permite comparar duas políticas sobre exatamente o mesmo conjunto de
tarefas, sem interferir em uma operação real \cite{law2000,silva2017}. O trabalho seguiu
as etapas abaixo.

\begin{alineas}
	\item definição do problema e das suposições do modelo;
	\item modelagem do layout genérico como grafo;
	\item definição do índice de prioridade e da regra de precedência;
	\item implementação do algoritmo de decisão e do simulador;
	\item verificação do funcionamento por testes automatizados;
	\item execução de cenários e comparação das duas políticas.
\end{alineas}

\subsection{Suposições do modelo}

As suposições adotadas na primeira versão são as seguintes. Elas ainda precisam
de validação em uma operação real.

\begin{alineas}
	\item cada rua é um beco sem saída, com acesso apenas pela frente; a travessa
	de fundo não está no grafo;
	\item os dois lados da rua contam como o mesmo trecho para efeito de
	bloqueio;
	\item apenas a empilhadeira \emph{parada operando} (posicionando, elevando e
	guardando) bloqueia o trecho; quem está em movimento passa;
	\item cada operador transporta um palete por viagem e retorna ao recebimento
	depois de guardá-lo;
	\item o endereço de cada material é fixo; categoria define a rua e curva
	define o nível;
	\item os tempos de percurso e de operação são determinísticos.
\end{alineas}

\subsection{Dados de entrada}

O catálogo de materiais tem 24 itens, 6 por rua, sendo dois de cada curva (A, B
e C). A rua depende da categoria do material, o nível depende da curva e a
posição e o lado dentro da rua são sorteados uma única vez com semente fixa,
de modo que o endereço de um material nunca muda entre execuções.

As tarefas de armazenagem são geradas a partir desse catálogo. Para cada tarefa,
sorteia-se a rua com probabilidades 0,5, 0,2, 0,2 e 0,1 para as ruas R1 a R4,
escolhe-se um material da rua com igual probabilidade e sorteia-se o instante
de criação entre 0 e 60~s. A concentração na primeira rua é intencional, pois é
nela que o congestionamento aparece. Cada cenário é identificado por uma
semente, e o gerador pseudoaleatório com semente garante que as duas políticas
recebam exatamente as mesmas tarefas \cite{law2000}.

\subsection{Modelo do grafo}

O layout é representado por um grafo não direcionado com 38 nós e 37 arestas. Os
nós são o ponto de recebimento, os cinco cruzamentos da travessa frontal e as
32 posições de armazenagem (8 por rua). As arestas são trechos de corredor,
com distância e tempo de percurso (velocidade de 2~m/s). Cada rua é um ramo que
sai de um único cruzamento de entrada e termina na posição mais funda.

A rota de cada operador é calculada pelo algoritmo de \citeonline{dijkstra1959}.
O custo de entrar em um nó é o tempo da aresta somado ao tempo restante da
empilhadeira que estiver parada operando naquele nó, o que torna o peso
dinâmico. A implementação percorre linearmente os nós ainda abertos para
escolher o próximo, com custo proporcional ao quadrado do número de nós, o que
é desprezível para 38 nós. Como o grafo recalcula a rota a cada decisão, ele
acompanha as mudanças de ocupação ao longo da simulação, ao contrário de uma
solução estática \cite{ramalingam1996}.

Uma observação importante: com a travessa de fundo fora do grafo, ele é uma
árvore, e existe um único caminho entre quaisquer dois nós. Por isso, neste
layout, o peso dinâmico não altera \emph{qual} caminho é percorrido; ele altera
o custo previsto do caminho, que entra na prioridade da tarefa. Rotas
alternativas só passam a existir se o layout tiver ciclos, por exemplo com a
travessa de fundo.

\subsection{Tempos de operação}

O tempo em que a empilhadeira permanece parada, ocupando o trecho, é dado por
\begin{equation}
	t_{op}(i) = t_0 + t_n \cdot n(i),
\end{equation}
em que $t_0$ é o tempo base de posicionar e guardar o palete, $t_n$ é o tempo
de elevação por nível e $n(i)$ é o nível do endereço da tarefa $i$. O tempo de
presença da tarefa dentro da rua inclui o deslocamento de ida e volta até a
posição $p(i)$:
\begin{equation}
	t_{rua}(i) = \frac{2 \, p(i) \, d}{v} + t_{op}(i),
\end{equation}
sendo $d$ a distância entre posições consecutivas e $v$ a velocidade. O valor
máximo $t_{rua}^{max}$ corresponde à posição mais funda no nível mais alto. Com
os parâmetros adotados (Tabela~\ref{tab:parametros}), $t_{rua}^{max} = 70$~s.

\subsection{Índice de prioridade}

Para cada tarefa $i$ que disputa a mesma rua, o índice de prioridade é
\begin{equation}
	P(i) = w_1 \, G(i) + w_2 \, Pos(i) - w_3 \, E(i),
\end{equation}
em que:
\begin{itemize}
	\item $G(i) \in [0,1]$ é o giro normalizado do material, obtido da curva
	ABC (A $=1$, B $=0{,}5$, C $=0{,}1$), seguindo a ideia de classificar os
	materiais por giro \cite{hausman1976};
	\item $Pos(i) = 1 - t_{rua}(i)/t_{rua}^{max}$ valoriza as tarefas que
	ocupam a rua por menos tempo, o que reflete a profundidade da posição e o
	nível de elevação;
	\item $E(i) = \min(1, e(i)/e_{ref})$ é a penalidade dinâmica pela espera
	prevista $e(i)$, calculada como a soma do tempo restante das empilhadeiras
	paradas no caminho da tarefa, com $e_{ref} = 40$~s como espera de
	referência;
	\item $w_1$, $w_2$ e $w_3$ são pesos parametrizáveis.
\end{itemize}

Há uma tensão a declarar: priorizar apenas o alto giro pode atrasar tarefas
rápidas, e priorizar apenas tarefas curtas pode deixar o alto giro esperando.
Os pesos controlam esse equilíbrio, e a simulação mostra o efeito de cada
escolha. Como a curva define o nível, $G(i)$ e $Pos(i)$ são correlacionados:
a curva A fica no nível mais baixo, com elevação curta, e a curva C fica no
nível mais alto, bloqueando a rua por mais tempo.

\subsection{Precedência por profundidade}

Em cada rua, somente a tarefa de destino mais fundo está liberada para
execução. Assim, quem armazena perto da entrada depois não bloqueia quem ainda
precisa chegar ao fundo. O índice $P(i)$ atua como critério de escolha entre as
tarefas liberadas de diferentes ruas e como desempate. Essa regra pode ser
desligada para fins de comparação, o que permite isolar o efeito dela do efeito
do giro.

\subsection{Algoritmo de decisão}

Sempre que um operador fica ocioso e existe tarefa na fila, o algoritmo executa
os passos a seguir. A decisão é tomada a cada evento, sobre o estado atual das
empilhadeiras.

\begin{alineas}
	\item monta o conjunto de tarefas liberadas, aplicando a regra de
	precedência por profundidade (ou toda a fila, se a regra estiver desligada);
	\item para cada tarefa liberada, calcula $G(i)$, $Pos(i)$ e a espera prevista
	$e(i)$ pelo estado atual das empilhadeiras que estão operando;
	\item calcula $P(i)$ e escolhe a tarefa de maior valor, mantendo a mais
	antiga em caso de empate;
	\item atribui a tarefa ao operador, que pega o palete no recebimento (5~s);
	\item calcula a rota pelo recebimento até o endereço, com o peso dinâmico
	de ocupação, e inicia o deslocamento;
	\item ao chegar, o operador fica parado operando durante $t_{op}(i)$, tempo em
	que bloqueia o trecho; depois retorna ao recebimento e fica ocioso.
\end{alineas}

A política de referência, ``ordem de chegada'', pula os passos 1 a 3 e sempre
escolhe a tarefa mais antiga da fila.

\subsection{Simulação}

A simulação avança por passos de tempo de 0,5~s, com limite de 20000~s por
execução, e termina quando todas as tarefas foram guardadas. Cada operador
passa pelos estados \emph{ocioso}, \emph{pegando}, \emph{indo},
\emph{operando} e \emph{voltando}. Quando o próximo nó do caminho está ocupado
por uma empilhadeira operando, o operador espera, e o tempo de espera é
acumulado por operador e por rua. Cada vez que um operador passa a esperar,
conta-se um bloqueio.

As duas políticas são executadas sobre cópias das mesmas tarefas, e a única
diferença entre as execuções é a regra de escolha da tarefa. As métricas
coletadas estão na Tabela~\ref{tab:metricas}.

\begin{table}[htb]
	\centering
	\caption{Métricas coletadas em cada simulação}
	\label{tab:metricas}
	\begin{tabular}{lp{8cm}}
		\toprule
		Métrica & Definição \\
		\midrule
		Tempo total & instante em que a última tarefa é guardada \\
		Espera por bloqueio & soma, entre os operadores, do tempo parado esperando um trecho ocupado \\
		Distância & soma das distâncias percorridas por todos os operadores \\
		Bloqueios & número de vezes que um operador passou a esperar \\
		Ciclo médio & média de (instante de guarda $-$ instante de criação) das tarefas \\
		Ciclo médio, curva A & mesmo cálculo, apenas para tarefas de curva A \\
		\bottomrule
	\end{tabular}
\end{table}

Os parâmetros adotados são apresentados na Tabela~\ref{tab:parametros}; todos
são configuráveis.

\begin{table}[htb]
	\centering
	\caption{Parâmetros do layout genérico e da simulação}
	\label{tab:parametros}
	\begin{tabular}{lr}
		\toprule
		Parâmetro & Valor \\
		\midrule
		Ruas de armazenagem & 4 \\
		Posições por rua & 8 \\
		Distância entre posições & 3~m \\
		Velocidade da empilhadeira & 2~m/s \\
		Tempo para pegar o palete & 5~s \\
		Tempo base de operação ($t_0$) & 10~s \\
		Tempo de elevação por nível ($t_n$) & 12~s \\
		Níveis por curva (A, B, C) & 1, 2, 3 \\
		Operadores & 5 \\
		Tarefas por cenário & 24 \\
		Passo de tempo & 0,5~s \\
		\bottomrule
	\end{tabular}
\end{table}

\subsection{Verificação e planejamento dos experimentos}

Para verificar se o simulador se comporta como o modelo descreve, foram escritos
sete testes automatizados, que conferem o endereço fixo por categoria e curva,
a rota até o fundo da rua, o efeito do peso dinâmico, a regra de precedência e
um cenário mínimo em que a ordem frente-antes-do-fundo gera bloqueio e a ordem
fundo-antes-da-frente não gera. A verificação confere se o modelo foi
implementado corretamente, e é diferente de validar se ele representa uma
operação real \cite{law2000}; esta segunda etapa não foi feita.

Como os resultados dependem das tarefas sorteadas, cada configuração é
executada em 30 cenários, com sementes de 1 a 30, e reporta-se a média. Os
experimentos comparam três configurações: ordem de chegada, algoritmo com os
pesos padrão ($w_1 = w_2 = w_3 = 0{,}5$) e regra de precedência, e algoritmo
sem a regra de precedência. Variar os pesos e o número de operadores fica para
a continuação do trabalho.

\subsection{Demonstração}

Como apoio, o mesmo algoritmo e a mesma simulação foram ligados a uma
visualização 3D e a um painel de indicadores. Essas interfaces apenas exibem o
resultado e não fazem parte da contribuição do trabalho.

% ----------------------------------------------------------
\section{Resultados preliminares}
\label{sec:resultados}
% ----------------------------------------------------------

Os resultados abaixo correspondem à média de 30 cenários, com 5 operadores e 24
tarefas por cenário, concentradas principalmente na primeira rua. Os valores
foram obtidos por simulação em layout genérico, com a primeira versão do
simulador, e não com dados de uma operação real.

\begin{table}[htb]
	\centering
	\caption{Comparação entre as políticas (média de 30 cenários)}
	\label{tab:resultados}
	\begin{tabular}{lccc}
		\toprule
		Indicador & Ordem de chegada & Algoritmo & Sem precedência \\
		\midrule
		Espera por bloqueio (s) & 318 & 399 & 324 \\
		Ciclo médio, curva A (s) & 193 & 157 & 98 \\
		Tempo total (s) & 430 & --- & 446 \\
		\bottomrule
	\end{tabular}
\end{table}

Com os pesos padrão ($w_1 = w_2 = w_3 = 0{,}5$) e a regra de precedência, o
algoritmo \emph{aumentou} a espera por bloqueio (399~s contra 318~s na ordem de
chegada). Em contrapartida, o tempo médio de ciclo da curva A caiu de 193~s para
157~s, e o ciclo médio geral ficou praticamente igual. Ao desligar a
precedência e priorizar apenas pelo índice, o ciclo da curva A caiu para 98~s e
a espera ficou em 324~s, equivalente à ordem de chegada, porém o tempo total
subiu de 430~s para 446~s.

Esses números indicam que a regra ``fundo primeiro'', isoladamente, não reduz a
espera total: a empilhadeira que armazena no fundo fica presa na saída enquanto
a da frente ainda está operando. O ganho observado vem da priorização do giro,
que reduz o tempo de ciclo dos materiais de curva A.

% ----------------------------------------------------------
\section{Considerações finais}
\label{sec:consideracoes}
% ----------------------------------------------------------

Este trabalho apresentou um algoritmo de priorização e roteirização de
empilhadeiras baseado em grafos dinâmicos e uma simulação para compará-lo com a
ordem de chegada. Os resultados preliminares
mostram que o critério de giro melhora o ciclo dos materiais de maior prioridade,
mas que a precedência por profundidade, como foi modelada, não basta para reduzir
os bloqueios.

\subsection{Limitações}

Este trabalho é um protótipo e não houve aplicação no mundo real. O layout
genérico mostra que o método é coerente, mas não mede o ganho real de um
centro logístico específico. Além disso, o modelo adota suposições que ainda
precisam de validação: ruas sem saída pelo fundo (a travessa de fundo não está
no grafo); os dois lados da rua contam como o mesmo trecho para efeito de
bloqueio; apenas a empilhadeira parada bloqueia o trecho; um palete por viagem;
e tempos fixos de operação. Como o grafo do layout é uma árvore, o peso dinâmico
não permite escolher rotas alternativas neste protótipo.

\subsection{Trabalhos futuros}

O próximo passo é fazer o algoritmo prever o conflito de \emph{saída}: a tarefa
da frente só começaria se terminasse antes de a empilhadeira do fundo precisar
sair. Em seguida, deve-se reavaliar a regra de precedência. Também ficam como
trabalho futuro a variação dos pesos e do número de operadores, a inclusão de
rotas alternativas no layout, a confirmação do tipo de estrutura de armazenagem
(porta-paletes seletivo ou estrutura profunda) e a aplicação com dados reais de
um centro logístico.

% ----------------------------------------------------------
% REFERÊNCIAS
% ----------------------------------------------------------
\postextual
\bibliography{referencias}

\end{document}
